%*******************************************************
% Acknowledgments
%*******************************************************
\pdfbookmark[1]{Acknowledgments}{acknowledgments}

\begin{flushright}{\slshape    
    We have seen that computer programming is an art, \\ 
    because it applies accumulated knowledge to the world, \\ 
    because it requires skill and ingenuity, and especially \\
    because it produces objects of beauty.} \\ \medskip
    --- \defcitealias{knuth:1974}{Donald E. Knuth}\citetalias{knuth:1974} \citep{knuth:1974}
\end{flushright}




\bigskip

\begingroup
\let\clearpage\relax
\let\cleardoublepage\relax
\let\cleardoublepage\relax
\chapter*{Acknowledgments}
This dissertation marks the end of my PhD journey, a journey that started long ago. Since I was a child, I have been fascinated by computers and their potential to solve complex problems. 
I always wanted to understand how they work and how to make them work better. This curiosity me to study programming and computer science since a young age.
When I was at 9th grade, I started attending high school computer science lessons, thaught by Tamir Moav and Miki TBF. They introduced me to the world of algorithms and programming, and I was liked. I spent countless hours learning and experimenting with different different programs and algorithms, and I was amazed by the power and elegance of computer science.

This led me to pursue a degree in computer science at the Open University of Israel, a place that allowed me to explore my interests and develop my skills, even at a young age. I was fortunate to have many inspiring professors and mentors who guided me along the way, and I am grateful for their support and encouragement.

After completing my undergraduate studies, I decided to go for a short trip to the industry in the world of software engineering. I worked as a software engineer for a year, and I learned a lot about the practical aspects of programming and software development. However, I realized that my true passion was in research and academia, and I decided to pursue a PhD in computer science. Nonetheless, I am grateful for the experience and the skills I gained during my time in the industry, as it helped me to become a better researcher and a more well-rounded computer scientist. For this, I would like to thank my colleagues and mentors in the industry, who taught me valuable lessons about teamwork, communication, and problem-solving. Specifically, I would like to thank my former colleagues Harel Fuchs, Shahar Lupu, Yossi Babayof and Zahi Abou.

Following my the quick stint in the industry, I moved to Leiden for my master's degree in computer science, where I was introduced to the world of research and academia. The program at LIACS greatly expanded my knowledge and skills in computer science, with a focus on research. I would like to thank the professors and PhD students who guided me, specifically, Jan van Rijn, Wojtek Kowalczyk, E
\endgroup



